\chapter{Informes, modelos, tesorería y recurrentes}\label{cap:informes}

\begin{objbox}
Verás cómo va el negocio por trimestre, sacarás las cifras de los modelos fiscales, estimarás la tesorería de los próximos meses y automatizarás facturas y gastos que se repiten.
\end{objbox}

\section{Informes}
\figura[0.6]{e21-informes}{Informes por año y trimestre.}
\begin{leyenda}
\item[\mk{1}] \textbf{PDF}. Descarga el informe del periodo.
\item[\mk{2}] \textbf{Año}.
\item[\mk{3}] \textbf{Periodo}: año completo o un trimestre (1T a 4T).
\item[\mk{4}] \textbf{Exportar}: \emph{Todo en Excel} (facturas, gastos y clientes en hojas separadas) o cada tabla en CSV. Es lo que normalmente pide una gestoría.
\end{leyenda}

\section{Modelos fiscales}
\figura[0.6]{e22-modelos}{Cifras orientativas de los modelos 130, 303, 115 y 347.}
\begin{leyenda}
\item[\mk{1}] \textbf{Año}.
\item[\mk{2}] \textbf{Trimestre}.
\item[\mk{3}] \textbf{CSV}. Descarga las cifras de cada modelo.
\end{leyenda}
REM calcula, con lo registrado, los datos de partida de los modelos \textbf{130} (pago fraccionado de IRPF, acumulado), \textbf{303} (IVA), \textbf{115} (retenciones por alquileres) y \textbf{347} (operaciones con terceros por encima de 3\,005,06~€).

\begin{warnbox}
Son cifras \textbf{orientativas}. REM no presenta nada ante Hacienda: pueden faltar deducciones, regularizaciones o pagos hechos fuera de la aplicación. Revísalas con tu asesor antes de presentar.
\end{warnbox}

\section{Tesorería}
\figura[0.6]{e23-tesoreria}{Previsión de tesorería a 30, 60 o 90 días.}
\begin{leyenda}
\item[\mk{1}] \textbf{Horizonte}: 30, 60 o 90 días.
\item[\mk{2}] \textbf{Opciones}: incluir facturas ya vencidas y estimaciones de gastos.
\end{leyenda}
La previsión usa tus vencimientos reales, el \textbf{retraso típico con que paga cada cliente} (calculado con tus facturas ya cobradas), tus recurrentes y tu histórico de gastos. Si introduces el saldo actual del banco, muestra también el saldo previsto y el \emph{punto más bajo}, para anticipar si te quedarás corto. Es una \textbf{estimación}, no una promesa.

\section{Recurrentes}
\figura[0.55]{e24-recurrentes}{Facturas y gastos recurrentes.}
\begin{leyenda}
\item[\mk{1}] \textbf{Preparar un borrador ahora}, sin cambiar la fecha de la siguiente.
\item[\mk{2}] \textbf{Editar}.
\item[\mk{3}] \textbf{Eliminar}.
\item[\mk{4}] \textbf{Nuevo gasto recurrente}.
\end{leyenda}
Una factura recurrente se crea desde una factura existente (\ui{Repetir…}) con una frecuencia (mensual, trimestral o anual). REM prepara el \textbf{borrador} en la fecha prevista; no lo emite ni lo envía. Un gasto recurrente, en cambio, se registra solo en cada periodo (uno por cada periodo que faltó si la aplicación estuvo cerrada).
